% ============================================================
%  插图与表格示例
%  不喜欢这页的话，把 main.tex 里的 % ============================================================
%  插图与表格示例
%  不喜欢这页的话，把 main.tex 里的 % ============================================================
%  插图与表格示例
%  不喜欢这页的话，把 main.tex 里的 % ============================================================
%  插图与表格示例
%  不喜欢这页的话，把 main.tex 里的 \input{chapters/figdemo} 删掉即可
% ============================================================

\section{插图与表格的写法}

\subsection{方法一：插入外部图片}

把图片文件放进 \texttt{images} 目录，正文里只写文件名，不写路径。XeLaTeX 下推荐 \texttt{PDF}（矢量，最清晰）和 \texttt{PNG}（位图，可透明），\texttt{JPG} 也可用，不建议 \texttt{EPS}。

\fig[0.62]{demo-parabola}{二次函数 $y=x^{2}-2x-3$ 的图象}{fig:parabola}

由图 \ref{fig:parabola} 可见，抛物线开口向上，顶点为 $(1,-4)$，与 $x$ 轴交于 $(-1,0)$ 和 $(3,0)$．

\begin{method}[图片宽度的控制]
\texttt{\textbackslash fig} 的第一个可选参数是相对版心宽度的比例：\texttt{[0.5]} 表示占半行宽，\texttt{[1.0]} 表示占满版心．文字较多的页面建议不超过 \texttt{[0.8]}．
\end{method}

\subsection{方法二：用 TikZ 直接画（数学图形首选）}

几何图形、函数图象、韦恩图这类示意图，用 TikZ 画比截图更清晰，缩放不失真，颜色和字号也能和正文统一．

\begin{figure}[H]
  \centering
  \begin{tikzpicture}[scale=1.15, >=Stealth]
    \draw[->] (-2.3,0) -- (2.5,0) node[below] {$x$};
    \draw[->] (0,-0.5) -- (0,2.6) node[left] {$y$};
    \draw[thick, mainblue, domain=-1.55:1.55, samples=100]
      plot (\x, {\x*\x}) node[right] {$y=x^{2}$};
    \node[below left] at (0,0) {$O$};
  \end{tikzpicture}
  \caption{用 TikZ 绘制的 $y=x^{2}$ 图象}
  \label{fig:tikz}
\end{figure}

\subsection{方法三：并排两张图}

\figpair{demo-parabola}{$y=x^{2}-2x-3$}{demo-parabola}{对比用同图示例}{并排插图示例}{fig:pair}

\subsection{表格}

\begin{table}[H]
  \centering
  \caption{集合与元素的关系}
  \label{tab:demo}
  \begin{tabular}{ccc}
    \toprule
    $a$ & $b$ & $a+b$ \\
    \midrule
    1 & 2 & 3 \\
    4 & 5 & 9 \\
    \bottomrule
  \end{tabular}
\end{table}

表 \ref{tab:demo} 中的三线表由 \texttt{booktabs} 宏包的 \verb|\toprule|、\verb|\midrule|、\verb|\bottomrule| 构成，比默认的表格线更适合教材排版．
 删掉即可
% ============================================================

\section{插图与表格的写法}

\subsection{方法一：插入外部图片}

把图片文件放进 \texttt{images} 目录，正文里只写文件名，不写路径。XeLaTeX 下推荐 \texttt{PDF}（矢量，最清晰）和 \texttt{PNG}（位图，可透明），\texttt{JPG} 也可用，不建议 \texttt{EPS}。

\fig[0.62]{demo-parabola}{二次函数 $y=x^{2}-2x-3$ 的图象}{fig:parabola}

由图 \ref{fig:parabola} 可见，抛物线开口向上，顶点为 $(1,-4)$，与 $x$ 轴交于 $(-1,0)$ 和 $(3,0)$．

\begin{method}[图片宽度的控制]
\texttt{\textbackslash fig} 的第一个可选参数是相对版心宽度的比例：\texttt{[0.5]} 表示占半行宽，\texttt{[1.0]} 表示占满版心．文字较多的页面建议不超过 \texttt{[0.8]}．
\end{method}

\subsection{方法二：用 TikZ 直接画（数学图形首选）}

几何图形、函数图象、韦恩图这类示意图，用 TikZ 画比截图更清晰，缩放不失真，颜色和字号也能和正文统一．

\begin{figure}[H]
  \centering
  \begin{tikzpicture}[scale=1.15, >=Stealth]
    \draw[->] (-2.3,0) -- (2.5,0) node[below] {$x$};
    \draw[->] (0,-0.5) -- (0,2.6) node[left] {$y$};
    \draw[thick, mainblue, domain=-1.55:1.55, samples=100]
      plot (\x, {\x*\x}) node[right] {$y=x^{2}$};
    \node[below left] at (0,0) {$O$};
  \end{tikzpicture}
  \caption{用 TikZ 绘制的 $y=x^{2}$ 图象}
  \label{fig:tikz}
\end{figure}

\subsection{方法三：并排两张图}

\figpair{demo-parabola}{$y=x^{2}-2x-3$}{demo-parabola}{对比用同图示例}{并排插图示例}{fig:pair}

\subsection{表格}

\begin{table}[H]
  \centering
  \caption{集合与元素的关系}
  \label{tab:demo}
  \begin{tabular}{ccc}
    \toprule
    $a$ & $b$ & $a+b$ \\
    \midrule
    1 & 2 & 3 \\
    4 & 5 & 9 \\
    \bottomrule
  \end{tabular}
\end{table}

表 \ref{tab:demo} 中的三线表由 \texttt{booktabs} 宏包的 \verb|\toprule|、\verb|\midrule|、\verb|\bottomrule| 构成，比默认的表格线更适合教材排版．
 删掉即可
% ============================================================

\section{插图与表格的写法}

\subsection{方法一：插入外部图片}

把图片文件放进 \texttt{images} 目录，正文里只写文件名，不写路径。XeLaTeX 下推荐 \texttt{PDF}（矢量，最清晰）和 \texttt{PNG}（位图，可透明），\texttt{JPG} 也可用，不建议 \texttt{EPS}。

\fig[0.62]{demo-parabola}{二次函数 $y=x^{2}-2x-3$ 的图象}{fig:parabola}

由图 \ref{fig:parabola} 可见，抛物线开口向上，顶点为 $(1,-4)$，与 $x$ 轴交于 $(-1,0)$ 和 $(3,0)$．

\begin{method}[图片宽度的控制]
\texttt{\textbackslash fig} 的第一个可选参数是相对版心宽度的比例：\texttt{[0.5]} 表示占半行宽，\texttt{[1.0]} 表示占满版心．文字较多的页面建议不超过 \texttt{[0.8]}．
\end{method}

\subsection{方法二：用 TikZ 直接画（数学图形首选）}

几何图形、函数图象、韦恩图这类示意图，用 TikZ 画比截图更清晰，缩放不失真，颜色和字号也能和正文统一．

\begin{figure}[H]
  \centering
  \begin{tikzpicture}[scale=1.15, >=Stealth]
    \draw[->] (-2.3,0) -- (2.5,0) node[below] {$x$};
    \draw[->] (0,-0.5) -- (0,2.6) node[left] {$y$};
    \draw[thick, mainblue, domain=-1.55:1.55, samples=100]
      plot (\x, {\x*\x}) node[right] {$y=x^{2}$};
    \node[below left] at (0,0) {$O$};
  \end{tikzpicture}
  \caption{用 TikZ 绘制的 $y=x^{2}$ 图象}
  \label{fig:tikz}
\end{figure}

\subsection{方法三：并排两张图}

\figpair{demo-parabola}{$y=x^{2}-2x-3$}{demo-parabola}{对比用同图示例}{并排插图示例}{fig:pair}

\subsection{表格}

\begin{table}[H]
  \centering
  \caption{集合与元素的关系}
  \label{tab:demo}
  \begin{tabular}{ccc}
    \toprule
    $a$ & $b$ & $a+b$ \\
    \midrule
    1 & 2 & 3 \\
    4 & 5 & 9 \\
    \bottomrule
  \end{tabular}
\end{table}

表 \ref{tab:demo} 中的三线表由 \texttt{booktabs} 宏包的 \verb|\toprule|、\verb|\midrule|、\verb|\bottomrule| 构成，比默认的表格线更适合教材排版．
 删掉即可
% ============================================================

\section{插图与表格的写法}

\subsection{方法一：插入外部图片}

把图片文件放进 \texttt{images} 目录，正文里只写文件名，不写路径。XeLaTeX 下推荐 \texttt{PDF}（矢量，最清晰）和 \texttt{PNG}（位图，可透明），\texttt{JPG} 也可用，不建议 \texttt{EPS}。

\fig[0.62]{demo-parabola}{二次函数 $y=x^{2}-2x-3$ 的图象}{fig:parabola}

由图 \ref{fig:parabola} 可见，抛物线开口向上，顶点为 $(1,-4)$，与 $x$ 轴交于 $(-1,0)$ 和 $(3,0)$．

\begin{method}[图片宽度的控制]
\texttt{\textbackslash fig} 的第一个可选参数是相对版心宽度的比例：\texttt{[0.5]} 表示占半行宽，\texttt{[1.0]} 表示占满版心．文字较多的页面建议不超过 \texttt{[0.8]}．
\end{method}

\subsection{方法二：用 TikZ 直接画（数学图形首选）}

几何图形、函数图象、韦恩图这类示意图，用 TikZ 画比截图更清晰，缩放不失真，颜色和字号也能和正文统一．

\begin{figure}[H]
  \centering
  \begin{tikzpicture}[scale=1.15, >=Stealth]
    \draw[->] (-2.3,0) -- (2.5,0) node[below] {$x$};
    \draw[->] (0,-0.5) -- (0,2.6) node[left] {$y$};
    \draw[thick, mainblue, domain=-1.55:1.55, samples=100]
      plot (\x, {\x*\x}) node[right] {$y=x^{2}$};
    \node[below left] at (0,0) {$O$};
  \end{tikzpicture}
  \caption{用 TikZ 绘制的 $y=x^{2}$ 图象}
  \label{fig:tikz}
\end{figure}

\subsection{方法三：并排两张图}

\figpair{demo-parabola}{$y=x^{2}-2x-3$}{demo-parabola}{对比用同图示例}{并排插图示例}{fig:pair}

\subsection{表格}

\begin{table}[H]
  \centering
  \caption{集合与元素的关系}
  \label{tab:demo}
  \begin{tabular}{ccc}
    \toprule
    $a$ & $b$ & $a+b$ \\
    \midrule
    1 & 2 & 3 \\
    4 & 5 & 9 \\
    \bottomrule
  \end{tabular}
\end{table}

表 \ref{tab:demo} 中的三线表由 \texttt{booktabs} 宏包的 \verb|\toprule|、\verb|\midrule|、\verb|\bottomrule| 构成，比默认的表格线更适合教材排版．
